% Self-contained universal upper bound; the final rounding implication is
% explicitly conditional until the goal global scalar certificate is accepted.
\begin{proposition}[An elementary universal upper bound]
\label{prop:new-self-contained-upper}
With $c=\log(1+\sqrt2)$,
\[
 G_{\mathbb R}\leq\frac{\pi}{2c}
                <\frac{1760}{987}<\frac{357}{200}=1.785.
\]
\end{proposition}
\begin{proof}
For a real unit vector $u$, in the Hilbert direct sum of its odd
real tensor powers define
\[
 \Phi(u)=\bigoplus_{k\geq0}
       \sqrt{\frac{c^{2k+1}}{(2k+1)!}}\,u^{\otimes(2k+1)},\qquad
 \Psi(u)=\bigoplus_{k\geq0}(-1)^k
       \sqrt{\frac{c^{2k+1}}{(2k+1)!}}\,u^{\otimes(2k+1)}.
\]
Since $\sinh c=1$, both vectors have norm one and
\[
 \langle\Phi(u),\Psi(v)\rangle
                    =\sin\bigl(c\langle u,v\rangle\bigr).
\]
Also $c<1$, because $\sinh1>1$, whereas
$\pi=4\int_0^1(1+x^2)^{-1}\,dx>2$; hence $c<\pi/2$.
For any finite families $u_i,v_j$, take a standard Gaussian vector
in the finite-dimensional span of their transformed vectors.  Put
$\epsilon_i=\operatorname{sgn}\langle G,\Phi(u_i)\rangle$ and
$\delta_j=\operatorname{sgn}\langle G,\Psi(v_j)\rangle$.
Two Gaussian linear forms with correlation $\rho$ have sign-product
expectation $2\arcsin(\rho)/\pi$: in their two-dimensional span,
rotational invariance makes the angle uniform and their signs
 differ with probability $\arccos(\rho)/\pi$.  Consequently
\[
 \mathbb E[\epsilon_i\delta_j]
 =\frac2\pi\arcsin\!\left(\sin(c\langle u_i,v_j\rangle)\right)
 =\frac{2c}{\pi}\langle u_i,v_j\rangle.
\]
For every real matrix $A$, its scalar sign norm bounds each rounded
bilinear form, so
\[
 \frac{2c}{\pi}\sum_{i,j}a_{ij}\langle u_i,v_j\rangle
 =\mathbb E\sum_{i,j}a_{ij}\epsilon_i\delta_j\leq C(A).
\]
Taking the vector supremum and then the supremum over matrices
proves $G_{\mathbb R}\leq\pi/(2c)$.

It remains to certify the numerical comparisons using rational
arithmetic.  The exact identity
\[
 \frac{x^4(1-x)^4}{1+x^2}
 =x^6-4x^5+5x^4-4x^2+4-\frac4{1+x^2}
\]
gives
\[
 0<\int_0^1\frac{x^4(1-x)^4}{1+x^2}\,dx
                         =\frac{22}{7}-\pi.
\]
Set $x_0=141/160$.  The positive exponential series and the
successive-tail ratio bound $x_0/(k+1)\leq x_0/14$ for $k\geq13$
give
\begin{align*}
 e^{x_0}
 &<\sum_{k=0}^{12}\frac{x_0^k}{k!}
       +\frac{x_0^{13}}{13!}\frac1{1-x_0/14}\\
 &=\frac{36546709859637513399380063330480378767}
          {15140013877549634840166400000000000000}
  <\frac{239}{99}<1+\sqrt2.
\end{align*}
The middle comparison follows by clearing positive denominators;
the final one follows from
$2-(140/99)^2=2/9801>0$.  Thus $c>141/160$, and
\[
 \frac\pi{2c}<\frac{22/7}{2(141/160)}
                =\frac{1760}{987},\qquad
 \frac{357}{200}-\frac{1760}{987}=\frac{359}{197400}>0.
\]
All polynomial and rational computations are exactly replayed by
\texttt{certify\_elementary\_upper.py}; they require no decimal
approximation to $\pi$, a logarithm, or a square root.
\end{proof}

\begin{proposition}[The three-significant-digit implication]
\label{prop:new-three-significant-digit-implication}
If a separately accepted complete matrix-sequence certificate proves
$G_{\mathbb R}\geq71/40$, then
\[
 \frac{71}{40}=1.775\leq G_{\mathbb R}
                         <\frac{1760}{987}<1.785.
\]
Thus $G_{\mathbb R}$ rounds to $1.78$ to three significant digits.
In particular, a complete scalar cap $N(T)\leq40/71$ for a seed
whose vector numerator tends to one proves this implication.
\end{proposition}
\begin{proof}
The upper endpoint is Proposition~\ref{prop:new-self-contained-upper}.
The interval $[1.775,1.785)$ is precisely the rounding interval for
$1.78$ at the hundredths place.  For the stated seed, its actual
matrix ratio limit is $1/N(T)\geq71/40$, so its matrices give the
required lower bound for the unrestricted constant.
\end{proof}
